% ============================================================
% APPENDIX D — Design of Experiments
% Pre-registered experimental protocol for the s035 model-ladder.
% Covers hypotheses, design matrix, controlled variables, statistical
% tests, kill rules, and change control.  Written before results.
% ============================================================
\section{Design of Experiments}
\label{app:doe}

This appendix documents the pre-registered experimental protocol for the
representation ladder (s035). The protocol was locked before Phase~1
execution and amended twice under change control (Table~\ref{app:tab:doe-versions}).
All hypotheses, statistical tests, and kill rules were specified before
any model was trained. Post-hoc adjustments to pass verification are
forbidden by the verification gates (Appendix~\ref{app:agent-lifecycle},
Table~\ref{app:tab:verify-lifecycle}).

\subsection{Hypotheses}
\label{app:doe-hypotheses}

Four hypotheses are tested across the representation ladder. H2a is the
single pre-registered confirmatory test; the others are descriptive or
exploratory.

\begin{table*}[t]
\centering
\caption{Pre-registered hypotheses. H2a is confirmatory (Wilcoxon + effect
  size); all others are descriptive or exploratory. Kill rules halt the
  experiment early if the baseline fails.}
\label{app:tab:hypotheses}
\small
\begin{tabular}{@{}p{0.8cm}p{4.5cm}p{3.5cm}p{2.5cm}p{3.0cm}@{}}
\toprule
\textbf{ID} & \textbf{Statement} & \textbf{Test} & \textbf{Pass criterion} & \textbf{Type} \\
\midrule
H1 &
  R0 features achieve skill above the mean predictor on $\geq 2$ of 3 targets &
  $R^2 > 0$ under spatial-blocked CV &
  $\geq 2$ targets with $R^2 > 0$ &
  Descriptive baseline gate \\[4pt]
H2a &
  R1 improves over R0 at the fold level &
  Wilcoxon signed-rank on paired (R0, R1) fold metrics &
  $p < 0.05$ AND Cohen's $d > 0.2$ &
  \textbf{Confirmatory} \\[4pt]
H2b &
  $\kappa_\text{geom}$ predicts which cells benefit more from R1 &
  Spearman $\rho$($\kappa_\text{geom}$, uplift) with bootstrap CI &
  $\rho > 0$, 95\% CI excludes zero &
  Exploratory ($n{=}9$) \\[4pt]
H3 &
  R2 temporal features improve over R1 when temporal mismatch is active &
  Paired delta: metric$_\text{R2}$ $-$ metric$_\text{R1}$ per fold &
  Uplift $> 3\%$ on primary target &
  Descriptive \\[4pt]
H4 &
  Audit flags predict representation uplift &
  Spearman $\rho$(flag count, uplift) &
  $\rho > 0.3$, bootstrap CI excludes zero &
  Exploratory \\
\bottomrule
\end{tabular}
\end{table*}

\subsection{Design matrix}
\label{app:doe-design}

The experiment is a controlled representation intervention: same folds,
same solver, same target---only the feature set changes. This is causal
at the pipeline level (representation is the sole intervention) but not
at the hydrology level (we do not randomize watersheds).

\begin{table}[t]
\centering
\caption{Design matrix. The only factor that changes across arms is the
  representation level (R0/R1/R2). All other factors are held constant
  or crossed.}
\label{app:tab:design-matrix}
\small
\begin{tabular}{@{}llc@{}}
\toprule
\textbf{Factor} & \textbf{Levels} & \textbf{Type} \\
\midrule
Scenario        & 5 metros (Table~\ref{app:tab:cells}) & Observational \\
Target          & 3 (claims, 311, HWM)                  & Fixed \\
Representation  & R0, R1, R2                            & \textbf{Intervention} \\
Solver          & HistGBDT, Ridge                       & Fixed \\
Split protocol  & random, spatial-blocked, leave-event-out & Fixed \\
Fold            & 5-fold spatial-blocked CV             & Fixed \\
\bottomrule
\end{tabular}
\end{table}

\paragraph{Modelable cells.}
Not all (scenario, target) pairs are modelable: \texttt{obs\_has\_311}
requires a 311 system (Houston, NYC only); \texttt{obs\_has\_hwm}
requires high-water marks (Houston, New~Orleans only). The nine modelable
cells are enumerated in Table~\ref{app:tab:cells}.

\subsection{Representation levels}
\label{app:doe-repr}

Each level is a strict superset of the prior, ensuring that any
performance change is attributable to the added features, not to
dropped features or re-tuning.

\begin{table}[t]
\centering
\caption{Representation ladder. Each level adds features to the prior;
  no features are removed or re-tuned.}
\label{app:tab:repr-levels}
\small
\begin{tabular}{@{}lp{4.5cm}c@{}}
\toprule
\textbf{Level} & \textbf{Added features} & \textbf{$\sim n_\text{feat}$} \\
\midrule
R0 & Static tabular: ACS demographics, SVI, FEMA flood zones,
     terrain (TWI, slope), HIFLD infrastructure,
     temporally-gated NFIP history & 33--36 \\
R1 & R0 + hydrology (NHDPlus catchments, levee proximity,
     sewershed type) + W-matrix spatial lags (8 cols) & 58--63 \\
R2 & R1 + temporal event dynamics (MRMS precipitation timing,
     tide gauge anomaly, HURDAT2 track features) & 67--72 \\
\bottomrule
\end{tabular}
\end{table}

\paragraph{Feature count ranges.}
Approximate counts vary by 1--3 features across scenarios due to
scenario-specific columns (e.g., HCFCD drainage district for Houston,
sewershed type for NYC). Runtime validation is authoritative; documented
counts are approximate.

\subsection{Controlled variables}
\label{app:doe-controlled}

The following are held constant across all representation levels to
isolate the representation intervention:

\begin{table}[t]
\centering
\caption{Controlled variables. No hyperparameter tuning is performed;
  defaults are fixed before execution.}
\label{app:tab:controlled}
\small
\begin{tabular}{@{}lll@{}}
\toprule
\textbf{Variable} & \textbf{Value} & \textbf{Rationale} \\
\midrule
HistGBDT max\_iter      & 200     & No tuning \\
HistGBDT max\_depth     & 6       & No tuning \\
HistGBDT learning\_rate & 0.1     & No tuning \\
Ridge alpha             & 1.0     & No tuning \\
Random seed             & 42      & Reproducibility \\
Imputation (Ridge)      & median  & No tuning \\
Scaling (Ridge)         & StandardScaler & No tuning \\
Fold assignments        & Fixed, saved to S3 & Shared across R0/R1/R2 \\
\bottomrule
\end{tabular}
\end{table}

\subsection{Statistical tests}
\label{app:doe-tests}

\paragraph{H2a (confirmatory).}
Fold-level Wilcoxon signed-rank test on paired deltas
(metric$_\text{R1}$ $-$ metric$_\text{R0}$) across all folds and cells.
The test requires \emph{both} $p < 0.05$ and Cohen's $d > 0.2$; a small
$p$-value without meaningful effect size does not pass.

\paragraph{H2b (exploratory).}
Spearman rank correlation between $\kappa_\text{geom}$ (pre-training
geometric compatibility, computed from problem structure alone) and
per-cell uplift (R1 $-$ R0). Bootstrap 95\% CI ($B = 1000$) must exclude
zero. At $n{=}9$ cells this is explicitly exploratory.

\paragraph{All other tests.}
H1 is a descriptive gate (does skill exist?). H3 and H4 are descriptive
with bootstrap confidence intervals. No multiplicity correction is
applied because only H2a is confirmatory; the others are reported as
exploratory effect sizes with uncertainty.

\subsection{Kill rules}
\label{app:doe-kill}

Kill rules halt the experiment early and prevent wasted compute:

\begin{enumerate}
\item \textbf{H1 FAIL on all 3 targets.} Data quality is too low for any
  modeling. Stop. Do not proceed to R1.
\item \textbf{H2a: Wilcoxon $p \geq 0.05$ or Cohen's $d < 0.2$.}
  Report as null result (R1 does not improve on R0). Proceed to R2 but
  flag as negative evidence.
\item \textbf{All uplift $< 1\%$.} Representation differences are noise;
  R0 is sufficient. Report and stop.
\end{enumerate}

\subsection{DO NOT constraints}
\label{app:doe-donot}

These constraints are binding for the duration of the experiment:
\begin{itemize}
\item Do NOT change fold assignments between representation levels.
\item Do NOT tune hyperparameters per representation level.
\item Do NOT run on scenarios without completed data assembly.
\item Do NOT use GPU instances (tabular models only).
\item Do NOT add features to R0 after the baseline is locked.
\item Do NOT impute targets (drop rows with NaN targets).
\end{itemize}

\subsection{Four-document experiment governance}
\label{app:doe-governance}

Each experiment is governed by four complementary documents with
non-overlapping scope:

\begin{table}[t]
\centering
\caption{Four-document governance stack. Each document answers a different
  question; together they form the complete experiment record.}
\label{app:tab:four-docs}
\small
\begin{tabular}{@{}lp{5.0cm}@{}}
\toprule
\textbf{Document} & \textbf{Question answered} \\
\midrule
DOE (this appendix) &
  \emph{Why} are we running this and \emph{what} do we expect to find?
  Hypotheses, design matrix, statistical tests, kill rules. \\[4pt]
Experiment Contract &
  \emph{What} inputs and outputs does each phase need?
  Packages, S3 artifacts, feature schemas, verification gate definitions. \\[4pt]
Checklist &
  \emph{Where} are we in execution?
  Checkboxes, dates, job names, commit hashes. \\[4pt]
Verification Gates &
  \emph{How} do we know results are correct?
  V0--V5 gates with commands, pass criteria, failure protocol. \\
\bottomrule
\end{tabular}
\end{table}

The DOE is pre-registration: it prevents post-hoc hypothesis invention.
The contract ensures data dependencies are met before launch. The
checklist tracks progress. The verification gates (Appendix~\ref{app:agent-lifecycle})
ensure correctness. Amendments to the DOE are permitted under change
control (Table~\ref{app:tab:doe-versions}) but must be recorded before
the affected phase executes.

\subsection{Change control}
\label{app:doe-changes}

\begin{table}[t]
\centering
\caption{DOE version history. Amendments are permitted before the affected
  phase executes; post-hoc amendments are forbidden.}
\label{app:tab:doe-versions}
\small
\begin{tabular}{@{}lll@{}}
\toprule
\textbf{Version} & \textbf{Date} & \textbf{Change} \\
\midrule
v1.0 & 2026-05-29 & Initial DOE (draft) \\
v1.1 & 2026-05-29 & Locked for execution \\
v1.2 & 2026-06-01 & R1/R2 feature bundles replaced with verified S3 sources;
                     kappa proxy diagnostics added \\
v1.7 & 2026-06-03 & H2 redesigned: fold-level Wilcoxon + Cohen's $d$ \\
v1.8 & 2026-06-03 & H2b added: $\kappa_\text{geom}$ independence test \\
\bottomrule
\end{tabular}
\end{table}

\paragraph{Paper-safe framing.}
This experiment is causal at the pipeline level: same folds, same solver,
same target---only the representation changes. It is \emph{not} causal at
the hydrology level (we do not randomize watersheds). The correct framing
is: ``Under controlled representation intervention, scenarios flagged by
audit B.1 showed $X$\% larger improvement when hydrologic features were
added.'' The incorrect framing---``adding hydrologic features causes better
flood prediction''---is confounded by feature quality, spatial coverage,
and data vintage.
